
\documentclass{article}
\usepackage{colortbl}
\usepackage{makecell}
\usepackage{multirow}
\usepackage{supertabular}

\begin{document}

\newcounter{utterance}

\centering \large Interaction Transcript for game `cladder', experiment `full\_v1.5\_default', episode 4889 with refineonly{-}Qwen3.5{-}9B.
\vspace{24pt}

{ \footnotesize  \setcounter{utterance}{1}
\setlength{\tabcolsep}{0pt}
\begin{supertabular}{c@{$\;$}|p{.15\linewidth}@{}p{.15\linewidth}p{.15\linewidth}p{.15\linewidth}p{.15\linewidth}p{.15\linewidth}}
    \# & \multicolumn{2}{c}{Player} && \multicolumn{2}{c}{Game Master} \\
    \hline

    \theutterance \stepcounter{utterance}  
    & & & \multicolumn{4}{p{0.6\linewidth}}{
        \cellcolor[rgb]{0.9,0.9,0.9}{
            \makecell[{{p{\linewidth}}}]{
                \texttt{\tiny{[P1$\langle$GM]}}
                \texttt{You are an expert in causal inference. The following question is not a typical commonsense query, but rather a meticulously designed question created by a professor specializing in causal inference, intended to assess the students' mastery of the course content.} \\
\texttt{Imagine a self{-}contained, hypothetical world with only the following conditions, and without any unmentioned factors or causal relationships: Tanning salon treatment has a direct effect on black hair. Going to the beach has a direct effect on black hair. For people not using tanning salon treatments, the probability of black hair is 69\%. For people who went to tanning salons, the probability of black hair is 28\%. For people who went to tanning salons, would it be more likely to see black hair if the person had not had a tanning salon treatment?} \\
\texttt{Start your answer with 'Yes' or 'No', followed by additional reasoning or evidence to support your explanation.} \\
            }
        }
    }
     \\ \\

    \theutterance \stepcounter{utterance}  
    & \multicolumn{4}{p{0.6\linewidth}}{
        \cellcolor[rgb]{0.561,0.737,0.561}{
            \makecell[{{p{\linewidth}}}]{
                \texttt{\tiny{[P1$\rangle$GM]}}
                \texttt{Yes} \\
\\ 
\texttt{**Reasoning:**} \\
\\ 
\texttt{To determine the answer, we must compare the probability of the outcome (black hair) under the actual treatment condition versus the counterfactual condition (no treatment) for the specific population in question (people who went to tanning salons).} \\
\\ 
\texttt{1.  **Identify the Probabilities:**} \\
\texttt{    *   The probability of having black hair for people who *did* go to tanning salons (observed outcome): \$P(H | T) = 28\textbackslash{}\%\$.} \\
\texttt{    *   The probability of having black hair for people who *did not* go to tanning salons (counterfactual outcome): \$P(H | \textbackslash{}neg T) = 69\textbackslash{}\%\$.} \\
\\ 
\texttt{2.  **Analyze the Comparison:**} \\
\texttt{    The question asks if it would be *more likely* to see black hair if the person had not had the treatment. This requires comparing the counterfactual probability (\$69\textbackslash{}\%\$) against the observed probability (\$28\textbackslash{}\%\$).} \\
\\ 
\texttt{3.  **Evaluate the Result:**} \\
\texttt{    Since \$69\textbackslash{}\% > 28\textbackslash{}\%\$, the likelihood of having black hair is significantly higher in the scenario where the tanning salon treatment was not received. The treatment acts as a suppressor of black hair in this hypothetical world.} \\
\\ 
\texttt{4.  **Conclusion:**} \\
\texttt{    Because the probability increases from 28\% to 69\% upon removing the treatment, the answer to whether it would be "more likely" is affirmative. The mention of "going to the beach" is irrelevant to this specific comparison as the necessary counterfactual probability is explicitly provided.} \\
\\ 
\texttt{**Final Answer:** Yes} \\
            }
        }
    }
    & & \\ \\

    \theutterance \stepcounter{utterance}  
    & & & \multicolumn{2}{p{0.3\linewidth}}{
        \cellcolor[rgb]{0.9,0.9,0.9}{
            \makecell[{{p{\linewidth}}}]{
                \texttt{\tiny{[GM$|$GM]}}
                \texttt{Yes} \\
            }
        }
    }
    & & \\ \\

    \theutterance \stepcounter{utterance}  
    & & & \multicolumn{2}{p{0.3\linewidth}}{
        \cellcolor[rgb]{0.9,0.9,0.9}{
            \makecell[{{p{\linewidth}}}]{
                \texttt{\tiny{[GM$|$GM]}}
                \texttt{yes} \\
            }
        }
    }
    & & \\ \\

    \theutterance \stepcounter{utterance}  
    & & & \multicolumn{2}{p{0.3\linewidth}}{
        \cellcolor[rgb]{0.9,0.9,0.9}{
            \makecell[{{p{\linewidth}}}]{
                \texttt{\tiny{[GM$|$GM]}}
                \texttt{game\_result = WIN} \\
            }
        }
    }
    & & \\ \\

\end{supertabular}
}

\end{document}
